LORAD comprises a fixed infrared beacon, a dual-photodiode receiver, and a
decoder. This section describes the hardware, the link protocol, the receiver
front end and detection logic, and the sequence model used for decoding.

\subsection{Hardware}\label{sec:hardware}

Both endpoints are Raspberry Pi Pico 2~W boards (RP2350) running MicroPython.
They are physically separate and share no timebase or synchronization channel,
so every timing relationship the decoder exploits is recovered from the optical
signal alone.

The beacon drives a TSAL6200 infrared emitter through a 2N2222 common-emitter
switch from GPIO~16. A second core runs an HTTP server
that sets the broadcast message and the rate multiplier at runtime; the
transmit loop reads both once per message and is otherwise independent of
network activity. Start pulses nominally 50~ms wide were \emph{received} with a
mean width of 50.18~ms and a standard deviation of 0.44~ms. That figure bounds
transmit jitter rather than measuring it, since it also contains detection
jitter and the receiver's timing quantization, but at 0.9\% of the symbol width
it establishes that the concurrent server does not perturb transmit timing
enough to matter.

The receiver carries two BPW34 photodiodes, each feeding an MCP6002
transimpedance stage with 100~k$\Omega$ feedback on a single 3.3~V rail, read
on ADC0 and ADC1 (GPIO~26 and GPIO~27). The feedback resistor was corrected
from 1~M$\Omega$ to 100~k$\Omega$ during development; all data reported here
were collected after that change, with the two channels matched to a gain ratio
of 1.15. The detectors are mounted close to coaxial rather than angled apart,
which is why the normalized difference between them carries no usable bearing
information at present (Section~\ref{conclusion}).

Two properties of the converter matter for interpreting the amplitude figures
in Section~\ref{results}. The RP2350 samples with a 12-bit successive
approximation converter, which MicroPython's \texttt{read\_u16} reports on a
16-bit scale; counts quoted out of 65{,}535 therefore carry 12 bits of real
resolution. The converter is also multiplexed across both inputs, so charge
from the previously selected channel biases each reading; the first conversion
after each switch is discarded, at the cost of doubling the conversions per
sample.

Detection runs as a MicroPython polling loop, and that loop's period sets the
receiver's timing resolution. Fitting the periodicity of the 90{,}870 measured
pulse widths gives a loop period of 262.4~$\mu$s, or 3.81~kiloiterations per
second, recovered independently at each of the five data rates and consistent
across them to within 0.5\%. This single number bounds the link. Symbol widths
are quantized at roughly a quarter of a millisecond, and the shortest symbol
spans 38 loop iterations at 1$\times$ but only 6.4 at 6$\times$ and 4.8 at
8$\times$, which is the mechanism behind the symbol loss reported in
Section~\ref{results} and the reason the useful rate range ends where it does.

\subsection{Link Protocol}

The beacon transmits baseband on-off keying. Information is carried in the
\emph{duration} of each emitted pulse rather than in its amplitude or in a
modulated carrier. Three pulse widths form the symbol alphabet: 10~ms denotes a
logical zero, 25~ms a logical one, and 50~ms a frame-start marker. Every pulse
is followed by a fixed 10~ms idle period.

Fig.~\ref{fig:protocol} shows one transmitted byte. Framing operates at byte
granularity: each character is sent as a start pulse followed by its eight data
bits, least-significant-bit first, giving nine pulses per byte. A message
terminates with a newline, after which the transmitter idles for 500~ms and
repeats, so a receiver may join the stream at any point and synchronize on the
next start pulse. That rest is fifty times the inter-symbol gap, making message
boundaries unambiguous even when individual symbols are lost.

\begin{figure}[t]
\centering
\includegraphics[width=\columnwidth]{fig_protocol.pdf}
\caption{One transmitted byte. Symbol identity is carried by pulse width, not
amplitude. Each byte is preceded by a 50~ms start pulse and followed by eight
data bits, least-significant-bit first; a 10~ms gap follows every pulse.}
\label{fig:protocol}
\end{figure}

At nominal timing a byte occupies approximately 280~ms, of which 60~ms is
framing overhead, giving a throughput of 28.6~bit/s. Two properties follow.
First, the per-byte start pulse and the gap that follows it consume 21\% of
airtime. Second, because the zero and one symbols differ in duration,
throughput is payload-dependent: an all-zero byte transmits in 220~ms against
340~ms for all-ones, a 55\% spread. This dependence is revisited in
Section~\ref{results}, where it accounts for a substantial part of the
difference between two experiments.

To characterize the link across operating points, all four timing constants are
divided by a common rate multiplier, preserving their ratios while compressing
the symbol set in time. The multiplier is set at runtime and applied at message
boundaries so that no transmission is split across two configurations. The
firmware clamps the multiplier to the range 0.5$\times$ to 8$\times$, nominally
14.3 to 228.6~bit/s, but only 1$\times$ through 6$\times$ were characterized.
The upper clamp is not a validated operating point: at 8$\times$ a zero symbol
spans 4.8 iterations of the 262.4~$\mu$s detection loop, which is too few to
time reliably.

\subsection{Receiver}

On the front end of Section~\ref{sec:hardware}, rather than thresholding
absolute signal level, which the DC bias of the photodiode network dominates,
the receiver applies a first-difference high-pass filter with leaky
integration,
\begin{equation}
y[n] = \bigl(x[n] - x[n-1]\bigr) + \alpha\, y[n-1], \qquad \alpha = 0.85,
\end{equation}
independently per channel, so that detection responds to signal transitions
rather than to absolute intensity. A rising edge is declared when the summed
filter output across both channels exceeds a threshold and a falling edge when
it falls below its negation; the interval between them is the symbol duration.

Edges are declared on the \emph{sum} of both channels rather than on either
individually. Off-axis, the further detector may barely register, and a
per-channel trigger would discard exactly those geometries.

The detection threshold adapts at runtime, tracking half the strongest recently
observed edge under an exponential moving average, clamped to a fixed range and
with a timeout that relaxes it by 10\% per second if no edge is detected. No
per-range tuning is required.

Serial output is buffered into preallocated arrays and emitted only during the
inter-message rest. An earlier implementation wrote each pulse as it was
detected, and because the sampling loop halts for the duration of a write, this
cost measurable symbols at higher rates, an effect quantified in
Section~\ref{experiments}. Reads of the reverse channel are confined to the
same idle window for the same reason.

\subsection{Decoder}

The decoder treats symbol recovery as sequence estimation. Hidden states are
positions within the byte framing; observations are detected pulses.
Fig.~\ref{fig:model} shows both halves of the model: the framing chain, which
is fixed by the protocol, and the three terms scored at each observation, of
which two carry fitted parameters.

\begin{figure*}[t]
\centering
\includegraphics[width=\textwidth]{fig_model.pdf}
\caption{The decoder. (a)~The framing state machine, fixed by the protocol and
not learned. (b)~The same model unrolled into a trellis. Columns are detected
pulses, rows are framing positions, grey edges are the legal advances, and the
bold path is the one Viterbi retains. An ordinary pulse advances one framing
position; the orange step at $o_3$ advances two on a single pulse, which is how
a symbol the receiver never detected appears. Only $\mu_s$, $\sigma_s^2$,
$\sigma_g$ and $\log\pi$ are fitted from data.}
\label{fig:model}
\end{figure*}

The model is a hidden Markov model and the recursion below is the ordinary
Viterbi algorithm, but three departures from the textbook formulation
$\lambda=(A,B,\pi_0)$ are worth naming, because each one simplifies or
complicates the standard picture in a way that matters here.

First, the transition structure $A$ is not a learned stochastic matrix. It is a
deterministic legality constraint read directly off the protocol: a transition
either exists or it does not, and Fig.~\ref{fig:model}(a) is the whole of it.
Nothing about the framing is estimated. Second, because the transmitted message
is known exactly, parameter fitting is fully supervised and closed-form. No
expectation-maximization over unlabeled sequences is required, so there are no
local optima and no convergence criterion; the estimator is the sample mean and
variance of a labeled set. Third, and most consequentially, a standard HMM
emits exactly one observation per transition. Here a single observation may
advance the chain by as many as $K+1$ transitions, because symbols the receiver
failed to detect still advance the framing even though they produced no pulse.
The closest standard analogues are a profile HMM with delete
states~\cite{Krogh1994} and the inner decoder of a watermark
code~\cite{Davey2001}, and it is this departure that Fig.~\ref{fig:model}(b) is
drawn to make visible: a skipped symbol is the one step in the trellis that
descends two rows in a single column.

Viterbi rather than the forward recursion is used because the required output
is a single self-consistent byte sequence, not a per-symbol posterior. A
marginal confidence at each position would not be constrained to spell an
admissible message.

\subsubsection{States and observations}

A state is the pair $q=(\text{bit position},\ \text{partial byte value})$,
giving $|Q| = 1 + \sum_{k=1}^{8} 2^{\,k-1} = 256$ states. Carrying the partial
byte in the state allows the payload constraint to prune during search: a path
that would complete an inadmissible byte is eliminated at the eighth data bit
rather than producing a candidate that must be rejected afterwards. The
admissible set is printable ASCII, 0x20 through 0x7E, together with the newline
the beacon appends as its terminator, which must remain admissible for the
final byte of a message to close.

The $t$-th detected pulse contributes the observation $o_t = (d_t,\ g_t)$,
where $d_t$ is its measured width and $g_t$ the interval since the previous
falling edge. Write $r$ for the rate multiplier, $w(s)$ for the nominal width
of symbol $s$ at $1\times$, and $G$ for the nominal inter-symbol gap.

\subsubsection{The three terms}

Each observation advances the chain by one emitting transition, optionally
preceded by up to $K=3$ \emph{skips} $u = (u_1,\dots,u_k)$ representing symbols
the receiver failed to detect. Three sources of evidence are combined in log
space.

The \emph{emission} term is a Gaussian class-conditional density over pulse
width,
\begin{equation}
\log p(d_t \mid s) = -\tfrac{1}{2}\left[\log\!\left(2\pi\sigma_s^2\right)
  + \frac{\left(d_t - \mu_s\right)^2}{\sigma_s^2}\right],
\label{eq:emission}
\end{equation}
with $\mu_s$ and $\sigma_s^2$ fitted per symbol and per rate. A generative form
is used because the Viterbi recursion requires $p(o \mid s)$ directly, rather
than a posterior that would have to be divided back through the class priors.

The \emph{gap} term evaluates the observed interval against the interval the
hypothesized skips predict. That prediction is structural, taken from nominal
protocol timing rather than fitted,
\begin{equation}
\hat{g}(k, u) = (k+1)\,\frac{G}{r} + \sum_{j=1}^{k} \frac{w(u_j)}{r},
\label{eq:gaphat}
\end{equation}
and only its spread is learned:
\begin{equation}
\log p(g_t \mid k, u) = \log \mathcal{N}\!\left(g_t;\ \hat{g}(k,u),\
  \bigl((1+k)\,\sigma_g\bigr)^2\right).
\label{eq:gap}
\end{equation}
The variance grows with $k$ because a hypothesis that stacks several
unobserved symbols accumulates their timing spread. Equation~\eqref{eq:gaphat}
is what makes a skipped symbol identifiable rather than merely permitted: its
value depends on \emph{which} symbols are hypothesized missing, not only how
many.

The \emph{skip prior} charges a fitted log-probability $\log\pi$ per unobserved
symbol. Without a real cost here the decoder could posit arbitrarily many
phantom symbols to force any byte into the admissible range.

\subsubsection{Recursion}

Let $\delta_t(q)$ be the best log-score of any path ending in state $q$ having
consumed $t$ observations, with $\delta_0$ zero at the start state and
$-\infty$ elsewhere. Then
\begin{equation}
\begin{split}
\delta_t(q') = \max_{q,\,u,\,s}\Bigl[\,&\delta_{t-1}(q) + k\log\pi \\
  &+ \gamma_t \log p(g_t \mid k,u) + \log p(d_t \mid s)\Bigr],
\end{split}
\label{eq:viterbi}
\end{equation}
maximized over source states $q$, skip sequences $u$ of length $k \le K$ that
walk $q$ to some intermediate state $m$ along legal transitions, and symbols
$s$ whose transition carries $m$ to $q'$. The gate $\gamma_t \in \{0,1\}$
suppresses the gap term where it does not apply: at $t=1$, because the first
pulse of a transmission is preceded by the 500~ms inter-message rest rather
than an inter-symbol gap, and wherever $g_t$ exceeds $1.5\,\hat{g}(K, \text{all
START})$, the largest interval the skip model could legitimately explain.
Decoding is the usual Viterbi backtrace from $\arg\max_q \delta_T(q)$, with
beam pruning that discards paths trailing the leader by more than a fixed
margin.

Fig.~\ref{fig:recovery} illustrates the mechanism
Equation~\eqref{eq:gaphat} encodes: an undetected symbol lengthens the
following gap by its own width plus one inter-symbol gap, so the observed gap
indicates both that a symbol was lost and which one.

\begin{figure}[t]
\centering
\includegraphics[width=\columnwidth]{fig_recovery.pdf}
\caption{Why an undetected symbol is recoverable. When the receiver misses a
pulse, the following inter-pulse gap is lengthened by that pulse's width plus
one gap period, so its duration identifies the lost symbol.}
\label{fig:recovery}
\end{figure}

\subsubsection{What is fitted, and how}

Eight scalars are learned per rate: three emission means, three emission
variances, one gap spread, and one skip prior. Everything else, the transition
structure, the nominal widths $w(\cdot)$ and $G$ entering
\eqref{eq:gaphat}, the skip bound $K$, and the beam width, is fixed by the
protocol or chosen once and held. Table~\ref{tab:params} lists the fitted
values.

Emission parameters are maximum-likelihood estimates over training pulses
labeled $s$,
\begin{equation}
\begin{split}
\mu_s &= \frac{1}{N_s}\sum_{i \in I_s} d_i, \\
\sigma_s^2 &= \max\!\left(\frac{1}{N_s}\sum_{i \in I_s} (d_i - \mu_s)^2,\
  (0.02\,\mu_s)^2\right).
\end{split}
\label{eq:fitem}
\end{equation}
The index set $I_s$ is drawn only from training transmissions in which
\emph{every} symbol was detected. That restriction is what makes the labels
exact: when the captured pulse count equals the transmitted symbol count, the
$i$-th pulse is the $i$-th symbol by position, so no alignment heuristic is
involved and the emission model cannot inherit an alignment error. The variance
floor prevents a class measured as near-deterministic from dominating every
path; it binds in 4 of the 15 cells, marked in Table~\ref{tab:params}.

The gap spread is the sample standard deviation of the departure of observed
gaps from nominal, over the same fully captured transmissions, floored the same
way:
\begin{equation}
\sigma_g = \max\!\left(\operatorname{sd}\bigl\{\,g_i - G/r\,\bigr\},\
  0.02\,G/r\right).
\label{eq:fitgap}
\end{equation}

The skip prior is the observed capture shortfall, Laplace-smoothed:
\begin{equation}
\log\pi = \log\frac{M + \tfrac{1}{2}}{E + 1},
\label{eq:fitskip}
\end{equation}
where $E$ is the total number of symbols transmitted across the training
transmissions and $M$ the number never detected. Smoothing is not cosmetic
here. At 28.6~bit/s the receiver drops nothing at all, so the unsmoothed
estimate is $\log 0$; the decoder must remain \emph{able} to posit a skip, just
very reluctant to. The smoothed value at that rate, $\log\pi = -10.08$, is
almost five nats more expensive than at 171.4~bit/s, which is the model
correctly learning that an unexplained symbol is far less plausible on a link
that is not dropping any.

\begin{table}[t]
\caption{Fitted parameters, per rate, from the training messages only. All
quantities in ms except $\log\pi$. Entries marked $\dagger$ are set by the
variance floor of \eqref{eq:fitem} rather than by the data.}
\label{tab:params}
\centering
\small
\setlength{\tabcolsep}{3.5pt}
\begin{tabular}{@{}lccccr@{}}
\toprule
Rate & \multicolumn{3}{c}{$\mu_s$ ($\sigma_s$)} & $\sigma_g$ & $\log\pi$ \\
\cmidrule(lr){2-4}
(bit/s) & ZERO & ONE & START & & \\
\midrule
28.6  & 10.18 (0.39) & 25.17 (0.50)$^\dagger$ & 50.17 (1.00)$^\dagger$ & 0.38 & $-10.08$ \\
57.1  & 5.17 (0.35)  & 12.65 (0.27)           & 25.15 (0.50)$^\dagger$ & 0.26 & $-5.80$ \\
85.7  & 3.50 (0.30)  & 8.47 (0.34)            & 16.82 (0.34)$^\dagger$ & 0.29 & $-5.50$ \\
114.3 & 2.66 (0.21)  & 6.38 (0.32)            & 12.63 (0.40)           & 0.32 & $-5.54$ \\
171.4 & 1.84 (0.22)  & 4.31 (0.15)            & 8.47 (0.31)            & 0.25 & $-5.19$ \\
\bottomrule
\end{tabular}
\end{table}

Table~\ref{tab:params} also answers a fair objection to the emission model:
if the protocol already specifies every symbol width, why learn $\mu_s$ at all?
Because the receiver does not measure the transmitted width. Every fitted mean
exceeds its nominal value, and by a strikingly stable amount: across all
fifteen rate-symbol combinations the excess is $155 \pm 16~\mu$s, independent of
both the symbol and the rate. That is 0.59 of the 262.4~$\mu$s detection loop
period, exactly the signature of a falling edge being registered on the first
loop iteration after it occurs. A decoder using nominal widths would place
every class centre roughly half a sample low; fitting absorbs the bias without
having to model its cause.

Because the decoder's lattice requires a backpointer per state per observation
(tens of thousands of entries for one message), it does not fit within the
receiver's memory, which raised allocation failures at 19~kB during development.
The receiver therefore streams per-pulse observations over its existing USB
serial connection and receives decoded messages back over the same link.
